\section{Допущения и ограничения}

\subsection{Допущения}

\begin{itemize}
  \item Основными пользователями являются студенты: они самостоятельно загружают материалы и получают отчёты с результатами проверки; при необходимости предусматриваются роли преподавателя и администратора.
  \item Документы для проверки передаются в систему в формате PDF (загрузка через веб-интерфейс или через интерфейс на основе электронной почты).
  \item Сервис itmo.id доступен для интеграции аутентификации; инфраструктура ИТМО предоставляет S3-совместимое хранилище и возможность развёртывания контейнеров и Nginx.
  \item Интерфейс ориентирован на русский язык.
  \item Используются типовые современные браузеры и подключение к сети.
\end{itemize}

\subsection{Ограничения}

\begin{itemize}
  \item Система выполняет только \emph{формальную} проверку документов; оценка содержания, качества изложения и научной корректности не входит в scope.
  \item Принимаются только документы в формате PDF; иные форматы (DOC, LaTeX-исходники и т.\,д.) не проверяются напрямую (проверяется итоговый PDF).
  \item Требования к оформлению и структуре соответствуют нормативным документам ИТМО и заданию по практике; изменения нормативов могут потребовать обновления правил проверки.
\end{itemize}
